\documentclass[11pt]{article}
\usepackage{mathblatt}


\begin{document}
\pfheftstil
\pfheftkopf{Lineare}{P10}
\pfrueckblick
\begin{pfaufg}{}{1.}{}
Gib an, auf welcher Achse der Punkt liegt: (0 | 2), ($-$2 | 0), (3 | 5).
\rechenraster{1}
\fusshilfe{\mbox{1:~(0 | 2) y-Achse; ($-$2 | 0) x-Achse; (3 | 5) auf keiner}}
\end{pfaufg}
\begin{pfaufg}{}{2.}{}
Ergänze die Tabelle Gleichung | m | n.
\par\smallskip\noindent \begingroup\renewcommand{\arraystretch}{1.3}\begin{tabular}{|c|c|c|}\hline \textbf{Gleichung} & \textbf{m} & \textbf{n} \\ \hline 3 | 1 & \rule{0pt}{3ex}\hspace{1.6cm} & \rule{0pt}{3ex}\hspace{1.6cm} \\ \hline $-$2 | 3 & \rule{0pt}{3ex}\hspace{1.6cm} & \rule{0pt}{3ex}\hspace{1.6cm} \\ \hline 4 | $-$2 & \rule{0pt}{3ex}\hspace{1.6cm} & \rule{0pt}{3ex}\hspace{1.6cm} \\ \hline 0,5 | $-$1 & \rule{0pt}{3ex}\hspace{1.6cm} & \rule{0pt}{3ex}\hspace{1.6cm} \\ \hline\end{tabular}\endgroup\par
\fusshilfe{\mbox{2:~3 | 1; $-$2 | 3; 4 | $-$2; 0,5 | $-$1}}
\end{pfaufg}
\begin{pfaufg}{}{3.}{}
Berechne: 2 $-$ ($-$1); 2 $-$ ($-$4); $-$1,5 $-$ 6; 3 $-$ ($-$2).
\rechenraster{1}
\fusshilfe{\mbox{3:~3; 6; $-$7,5; 5}}
\end{pfaufg}
\begin{pfaufg}{}{4.}{}
Berechne: $-$2 · ($-$4) + 2; 3 · 5 + 1; ½ · ($-$10) + 1; ($-$5)² $-$ 4.
\rechenraster{1}
\fusshilfe{\mbox{4:~10; 16; $-$4; 21}}
\end{pfaufg}
\begin{pfaufg}{}{5.}{}
Löse: 2x + 6 = 0; $-$0,5x + 10 = 0; 3x + 12 = 30.
\fusshilfe{\mbox{5:~x = $-$3; x = 20; x = 6}}
\end{pfaufg}
\pfstufekopf[serkennenundablesen]{erkennen und ablesen}{in 2 der letzten 5 Prüfungen}
\pfgruppe{}
\begin{pfaufg}{}{6.}{}
Unterstreiche in $f(x) = 3x - 5$ die Steigung m. Kreise den y-Achsenabschnitt n ein. Gib m und n an.
\par\noindent m = \leerfeld , n = \leerfeld \par
\fusshilfe{\mbox{6:~$-5$}}
\end{pfaufg}
\begin{pfaufg}{}{7.}{}
Kreuze an, ob die Gerade zu $f(x) = 3x - 4$ steigt oder fällt. 
\pfkreuzzeile{steigt}\pfkreuzzeile{fällt}
\fusshilfe{\mbox{7:~steigt}}
\end{pfaufg}
\begin{pfaufg}{}{8.}{}
Das Koordinatensystem zeigt eine Gerade mit einem Steigungsdreieck. Das Dreieck geht einen Schritt nach rechts. Lies ab, wie viele Schritte es nach oben oder unten geht.
\par\smallskip\noindent \begin{ksys}[xmin=-4,xmax=4,ymin=-5,ymax=5,ablesen] \gerade{3}{-1}{g} \steigungsdreieck{0}{-1}{3} \end{ksys}\par
\par\noindent \leerfeld[nach]  \leerfeld[\_\_] \par
\fusshilfe{\mbox{8:~$3$}}
\end{pfaufg}
\pfgruppe{}
\begin{pfaufg}{BY ’22}{9.}{}
Eine Bahnstrecke hat die Steigung 0,25. Ergänze: Bei einer waagrechten Entfernung von \_\_\_ m wird ein Höhenunterschied von \_\_\_ m überwunden.
\rechenraster{1}
\fusshilfe{\mbox{9:~z. B. 100 m waagrecht, 25 m Höhe}}
\end{pfaufg}
\begin{pfaufg}{P10 ’19}{10.}{2 BE}
In einem Koordinatensystem sind die Graphen $f$, $g$ und $h$ von drei linearen Funktionen gezeichnet. Schreib zu jeder Eigenschaft den Graphen ($f$, $g$ oder $h$), der sie hat: „Der Graph hat die Steigung $-2$.“ „Der Graph liegt parallel zur x-Achse.“
\par\smallskip\noindent \sachtabelle{lc}{Eigenschaft & Graph}{Steigung $-$2 & \leerzelle\\ parallel zur x-Achse & \leerzelle}\par
\fusshilfe{\mbox{10:~Anstieg $-$2: f; parallel zur x-Achse: h}}
\end{pfaufg}
\begin{pfaufg}{VERA ’14}{11.}{}
Gegeben sind die Geraden y = 4x + 3 und y = 4x $-$ 3. Wie liegen sie zueinander? Kreuze an: parallel / schneiden sich in einem Punkt / identisch. Begründe.
\rechenraster{1}
\end{pfaufg}
\pfgruppe{}
\begin{pfaufg}{}{12.}{}
Das Koordinatensystem zeigt die Gerade g. Lies ihren y-Achsenabschnitt n ab.
\par\smallskip\noindent \begin{ksys}[xmin=-5,xmax=5,ymin=-5,ymax=5,ablesen] \gerade{1}{3}{g} \end{ksys}\par
\par\noindent n = \leerfeld \par
\fusshilfe{\mbox{12:~$3$}}
\end{pfaufg}
\pfgruppe{}
\begin{pfaufg}{P10 ’16}{13.}{1 BE}
Im Koordinatensystem sind zwei Geraden $f$ und $g$ gezeichnet. Kreuze an, welche der beiden fällt:
\pfkreuzzeile{$f$}\pfkreuzzeile{$g$}
\par\smallskip\noindent \begin{tikzpicture}\begin{axis}[xmin=-1.5,xmax=3.5,ymin=-2.5,ymax=1.5,axis lines=middle,tick label style={font=\scriptsize},clip=true,/pgf/number format/use comma,axis line style={-{Stealth[length=2mm]}},every axis plot/.append style={line width=0.9pt},x=0.55cm,y=0.55cm,xtick distance=1,ytick distance=1,enlargelimits=false,grid=both,grid style={mbgitter,line width=0.3pt},major grid style={mbgitter},xlabel={$x$},ylabel={$y$},x label style={anchor=west},y label style={anchor=south}]\addplot[black,domain=-1.5:3.5,samples=120] {-1.5*x+0.5};\node[font=\small,fill=white,inner sep=1pt,anchor=south west] at (axis cs:1.78,-2.17) {$f$};\addplot[black,domain=-1.5:3.5,samples=120] {0.5*x-1.5};\node[font=\small,fill=white,inner sep=1pt,anchor=south west] at (axis cs:3.1,0.05) {$g$};\end{axis}\end{tikzpicture}\par
\fusshilfe{\mbox{13:~f}}
\end{pfaufg}
\pfgruppe{}
\begin{pfaufg}{P10 ’24}{14.}{1 BE}
Gegeben ist die Gerade $y = 0{,}5x - 1$. Zeichne den Punkt ein, an dem sie die y-Achse trifft.
\par\smallskip\noindent \begin{tikzpicture}\begin{axis}[xmin=-2,xmax=2,ymin=-2,ymax=2,axis lines=middle,tick label style={font=\scriptsize},clip=true,/pgf/number format/use comma,axis line style={-{Stealth[length=2mm]}},every axis plot/.append style={line width=0.9pt},x=0.55cm,y=0.55cm,xtick distance=1,ytick distance=1,enlargelimits=false,grid=both,grid style={mbgitter,line width=0.3pt},major grid style={mbgitter},xlabel={$x$},ylabel={$y$},x label style={anchor=west},y label style={anchor=south}]\end{axis}\end{tikzpicture}\par
\rechenraster{1}
\fusshilfe{\mbox{14:~Punkt (0|$-$1) markiert}}
\end{pfaufg}
\pfgruppe{}
\begin{pfaufg}{P10 ’25}{15.}{1 BE}
Kreuze das Bild an, das den Graphen einer linearen Funktion zeigt.
\par\smallskip\noindent \begin{tabular}[t]{@{}c@{}}{\small A}\enspace$\square$\\\begin{tikzpicture}\begin{axis}[width=3.9cm,height=3.4cm,scale only axis=false,axis lines=middle,xtick=\empty,ytick=\empty,xmin=-0.8,xmax=1.036,ymin=-0.6,ymax=1.128,clip=true,axis line style={-{Stealth[length=1.5mm]}}]\addplot[black,line width=0.9pt,domain=-0.8:0.9,samples=80] {-0.6*x+0.2};\end{axis}\end{tikzpicture}\end{tabular}\hspace{0.4cm}\begin{tabular}[t]{@{}c@{}}{\small B}\enspace$\square$\\\begin{tikzpicture}\begin{axis}[width=3.9cm,height=3.4cm,scale only axis=false,axis lines=middle,xtick=\empty,ytick=\empty,xmin=-0.8,xmax=1.036,ymin=-0.6,ymax=1.128,clip=true,axis line style={-{Stealth[length=1.5mm]}}]\addplot[black,line width=0.9pt,domain=0.08:0.85,samples=80] {0.0680/x};\end{axis}\end{tikzpicture}\end{tabular}\hspace{0.4cm}\begin{tabular}[t]{@{}c@{}}{\small C}\enspace$\square$\\\begin{tikzpicture}\begin{axis}[width=3.9cm,height=3.4cm,scale only axis=false,axis lines=middle,xtick=\empty,ytick=\empty,xmin=-0.8,xmax=1.036,ymin=-0.6,ymax=1.128,clip=true,axis line style={-{Stealth[length=1.5mm]}}]\addplot[black,line width=0.9pt,domain=-0.6:0.6,samples=80] {3.5*(x^2)-0.5};\end{axis}\end{tikzpicture}\end{tabular}\hspace{0.4cm}\begin{tabular}[t]{@{}c@{}}{\small D}\enspace$\square$\\\begin{tikzpicture}\begin{axis}[width=3.9cm,height=3.4cm,scale only axis=false,axis lines=middle,xtick=\empty,ytick=\empty,xmin=-0.8,xmax=1.036,ymin=-0.6,ymax=1.128,clip=true,axis line style={-{Stealth[length=1.5mm]}}]\addplot[black,line width=0.9pt] coordinates {(0.1,-0.35) (0.5,0.4) (0.1,0.8)};\end{axis}\end{tikzpicture}\end{tabular}\par
\fusshilfe{\mbox{15:~erste Abbildung (fallende Gerade)}}
\end{pfaufg}
\pfgruppe{}
\begin{pfaufg}{P10 ’19}{16.}{1 BE}
In einem Koordinatensystem sind die Graphen $f$, $g$ und $h$ von drei linearen Funktionen gezeichnet. Gib die Nullstelle der Funktion $g$ an.
\par\smallskip\noindent \begin{tikzpicture}\begin{axis}[xmin=-5,xmax=5,ymin=-5,ymax=5,axis lines=middle,tick label style={font=\scriptsize},clip=true,/pgf/number format/use comma,axis line style={-{Stealth[length=2mm]}},every axis plot/.append style={line width=0.9pt},x=0.55cm,y=0.55cm,xtick distance=1,ytick distance=1,enlargelimits=false,grid=both,grid style={mbgitter,line width=0.3pt},major grid style={mbgitter},xlabel={$x$},ylabel={$y$},x label style={anchor=west},y label style={anchor=south}]\addplot[black,domain=-5:5,samples=120] {-2*x+2};\node[font=\small,fill=white,inner sep=1pt,anchor=south west] at (axis cs:2.88,-3.76) {$f$};\addplot[black,domain=-5:5,samples=120] {x+2};\node[font=\small,fill=white,inner sep=1pt,anchor=south west] at (axis cs:1.78,3.78) {$g$};\addplot[black,domain=-5:5,samples=120] {-2};\node[font=\small,fill=white,inner sep=1pt,anchor=south west] at (axis cs:4.2,-2) {$h$};\fill (axis cs:0,-2) circle (1.6pt) node[above right,font=\small,inner sep=1pt] {$A$};\fill (axis cs:2,-2) circle (1.6pt) node[above right,font=\small,inner sep=1pt] {$B$};\fill (axis cs:0,2) circle (1.6pt) node[above right,font=\small,inner sep=1pt] {$C$};\end{axis}\end{tikzpicture}\par
\rechenraster{1}
\fusshilfe{\mbox{16:~x = $-$2}}
\end{pfaufg}
\pfgruppe{}
\begin{pfaufg}{P10 ’19}{17.}{3 BE}
In einem Koordinatensystem sind die Graphen $f$, $g$ und $h$ von drei linearen Funktionen gezeichnet. Ordne jedem Graphen $f$, $g$ und $h$ eine Gleichung aus dieser Liste zu: $y = -2x + 2$, $y = -\tfrac{1}{2}x + 2$, $y = -2$, $y = x + 2$, $y = -2x$, $y = x - 2$.
\par\smallskip\noindent \sachtabelle{lc}{Graph & Gleichung}{f & \leerzelle\\ g & \leerzelle\\ h & \leerzelle}\par
\pfzwfrage{Lies für jeden Graphen den Schnittpunkt mit der y-Achse ab.}
\fusshilfe{\mbox{17:~f: y = $-$2x + 2; g: y = x + 2; h: y = $-$2}}
\end{pfaufg}

\pfstufekopf[sausgleichungzeichnen]{aus Gleichung zeichnen}{in 4 der letzten 5 Prüfungen}
\pfgruppe{}
\begin{pfaufg}{}{18.}{}
Fülle die Wertetabelle zu $f(x) = 3x - 2$ aus.
\par\smallskip\noindent \wertetabelle{x}{f(x)}{-1,0,1,2}\par
\fusshilfe{\mbox{18:~$-5$, $-2$, $1$, $4$}}
\end{pfaufg}
\begin{pfaufg}{BY ’21}{19.}{}
Zeichne die Gerade y = $-$2x + 4 in ein Koordinatensystem.
\par\smallskip\noindent \begin{tikzpicture}\begin{axis}[xmin=-1,xmax=5,ymin=-1,ymax=5,axis lines=middle,tick label style={font=\scriptsize},clip=true,/pgf/number format/use comma,axis line style={-{Stealth[length=2mm]}},every axis plot/.append style={line width=0.9pt},x=0.55cm,y=0.55cm,xtick distance=1,ytick distance=1,enlargelimits=false,grid=both,grid style={mbgitter,line width=0.3pt},major grid style={mbgitter},xlabel={$x$},ylabel={$y$},x label style={anchor=west},y label style={anchor=south}]\end{axis}\end{tikzpicture}\par
\rechenraster{1}
\fusshilfe{\mbox{19:~Gerade durch (0; 4) und (2; 0)}}
\end{pfaufg}
\begin{pfaufg}{NI ’22}{20.}{}
Zeichne den Graphen von y = $\tfrac{2}{3}$x $-$ 2 in ein Koordinatensystem.
\rechenraster{1}
\fusshilfe{\mbox{20:~Gerade durch (0; $-$2) und (3; 0)}}
\end{pfaufg}
\begin{pfaufg}{P10 ’21}{21.}{2 BE}
Die lineare Funktion $f$ hat die Gleichung $y = 3x + 1$. Zeichne den Graphen von $f$ in das Koordinatensystem.
\par\smallskip\noindent \begin{tikzpicture}\begin{axis}[xmin=-4,xmax=4,ymin=-4,ymax=4,axis lines=middle,tick label style={font=\scriptsize},clip=true,/pgf/number format/use comma,axis line style={-{Stealth[length=2mm]}},every axis plot/.append style={line width=0.9pt},x=0.55cm,y=0.55cm,xtick distance=1,ytick distance=1,enlargelimits=false,grid=both,grid style={mbgitter,line width=0.3pt},major grid style={mbgitter},xlabel={$x$},ylabel={$y$},x label style={anchor=west},y label style={anchor=south}]\end{axis}\end{tikzpicture}\par
\rechenraster{1}
\fusshilfe{\mbox{21:~Gerade durch (0 | 1) und (1 | 4)}}
\end{pfaufg}
\begin{pfaufg}{P10 ’24}{22.}{2 BE}
Die lineare Funktion $f$ hat die Gleichung $f(x) = 4x + 1$. Zeichne ihren Graphen in das Koordinatensystem.
\par\smallskip\noindent \begin{tikzpicture}\begin{axis}[xmin=-4,xmax=4,ymin=-5,ymax=8,axis lines=middle,tick label style={font=\scriptsize},clip=true,/pgf/number format/use comma,axis line style={-{Stealth[length=2mm]}},every axis plot/.append style={line width=0.9pt},x=0.42cm,y=0.42cm,xtick distance=1,ytick distance=1,enlargelimits=false,grid=both,grid style={mbgitter,line width=0.3pt},major grid style={mbgitter},xlabel={$x$},ylabel={$y$},x label style={anchor=west},y label style={anchor=south}]\end{axis}\end{tikzpicture}\par
\rechenraster{1}
\fusshilfe{\mbox{22:~Gerade durch (0|1) und (1|5), steil steigend}}
\end{pfaufg}
\pfgruppe{}
\begin{pfaufg}{P10 ’22}{23.}{3 BE}
Gegeben ist die lineare Funktion $f(x) = -2x + 3$. Zeichne den Graphen von $f$ in das Koordinatensystem. Gib die Nullstelle von $f$ an.
\par\smallskip\noindent \begin{tikzpicture}\begin{axis}[xmin=-4,xmax=5,ymin=-4,ymax=7,axis lines=middle,tick label style={font=\scriptsize},clip=true,/pgf/number format/use comma,axis line style={-{Stealth[length=2mm]}},every axis plot/.append style={line width=0.9pt},x=0.55cm,y=0.55cm,xtick distance=1,ytick distance=1,enlargelimits=false,grid=both,grid style={mbgitter,line width=0.3pt},major grid style={mbgitter},xlabel={$x$},ylabel={$y$},x label style={anchor=west},y label style={anchor=south}]\end{axis}\end{tikzpicture}\par
\pfzwfrage{Löse die Gleichung $-2x + 3 = 0$.}
\rechenraster{2}
\fusshilfe{\mbox{23:~Gerade durch (0|3) und (1|1); x = 1,5}}
\end{pfaufg}
\pfsteckt{steckt auch in Nr. 70 (P10 ’23)}
\pfstufekopf[sdurchzweipunktegleic]{durch zwei Punkte, Gleichung ablesen}{in 1 der letzten 5 Prüfungen}
\pfgruppe{}
\begin{pfaufg}{nach BW ’24}{24.}{}
Eine lineare Funktion geht durch (0 | 1) und (1 | 4). Bestimme ihre Gleichung.
\rechenraster{1}
\fusshilfe{\mbox{24:~y = 3x + 1}}
\end{pfaufg}
\begin{pfaufg}{HH ’26}{25.}{}
Die Gerade g geht durch A(3 | 4) und B(2 | 1). Ermittle ihre Steigung.
\rechenraster{1}
\fusshilfe{\mbox{25:~m = (4 $-$ 1) : (3 $-$ 2) = 3}}
\end{pfaufg}
\begin{pfaufg}{NRW ’23}{26.}{}
Eine Gerade hat die Steigung m = 3 und geht durch P(2 | 6). Gib ihre Gleichung an.
\rechenraster{1}
\fusshilfe{\mbox{26:~y = 3x}}
\end{pfaufg}
\pfgruppe{}
\begin{pfaufg}{}{27.}{}
Trage die Punkte $A(-3|-1)$ und $B(2|4)$ in das Koordinatensystem ein. Zeichne die Gerade durch beide Punkte.
\par\smallskip\noindent \begin{ksys}[xmin=-4,xmax=4,ymin=-4,ymax=5] \end{ksys}\par
\rechenraster{1}
\fusshilfe{\mbox{27:~Gerade durch $A(-3|-1)$ und $B(2|4)$}}
\end{pfaufg}
\begin{pfaufg}{P10 ’17}{28.}{3 BE}
Eine Gerade $g$ geht durch die Punkte $K(-4\,|\,-1)$ und $L(2\,|\,2)$. Zeichne ein Koordinatensystem und darin die Gerade $g$. Zeige, dass $y = \tfrac{1}{2}x + 1$ eine Gleichung von $g$ ist.
\par\smallskip\noindent \begin{tikzpicture}[x=0.5cm,y=0.5cm]\draw[mbgitter,line width=0.3pt] (0,0) grid (28,15);\end{tikzpicture}\par
\pfzwfrage{Berechne die Steigung aus den Koordinaten von K und L.}
\pfzwfrage{Setze L in $y = mx + n$ ein und berechne $n$.}
\rechenraster{3}
\fusshilfe{\mbox{28:~Gerade durch K und L; m = ½, n = 1 $\Rightarrow$ y = ½x + 1}}
\end{pfaufg}
\begin{pfaufg}{P10 ’15}{\llap{\textasteriskcentered\,}29.}{3 BE}
Lena springt mit dem Fallschirm aus einem Flugzeug. Ihre Höhe $h$ (in Metern) hängt von der Zeit $t$ (in Sekunden) ab. Nach 20 Sekunden freiem Fall öffnet sie in 1200 m Höhe den Schirm; 100 Sekunden später ist sie noch 700 m hoch. Während sie mit offenem Schirm gleitet, liegt der Graph auf einer Geraden. Stelle eine Gleichung dieser Geraden auf.
\par\smallskip\noindent \begin{tikzpicture}\begin{axis}[xmin=0,xmax=280,ymin=0,ymax=2400,tick label style={font=\scriptsize},clip=true,/pgf/number format/use comma,axis line style={-{Stealth[length=2mm]}},every axis plot/.append style={line width=0.9pt},width=11.5cm,height=7.5cm,grid=both,grid style={mbgitter,line width=0.3pt},minor tick num=1,axis lines=left,xlabel={Zeit t in Sekunden},ylabel={Höhe h in Metern},label style={font=\small},scaled ticks=false,/pgf/number format/1000 sep={\,}]\addplot[black,domain=0:20,samples=120] {-3*(x^2)+2400};\node[font=\small,fill=white,inner sep=1pt,anchor=south west] at (axis cs:18.4,1384.32) {$G$};\addplot[black,domain=20:260,samples=120] {-5*x+1300};\end{axis}\end{tikzpicture}\par
\pfzwfrage{Berechne die Steigung aus den Punkten (20 | 1200) und (120 | 700).}
\rechenraster{3}
\fusshilfe{\mbox{29:~h(t) = $-$5t + 1300}}
\end{pfaufg}
\pfgruppe{}
\begin{pfaufg}{NRW ’21}{30.}{}
Gib eine lineare Gleichung an, die zur Wertetabelle passt.
\par\smallskip\noindent \sachtabelle{l}{x 0, 1, …}{y 2, 3,5, …}\par
\fusshilfe{\mbox{30:~m = (3,5 $-$ 2) : (1 $-$ 0) = 1,5; y = 1,5x + 2}}
\end{pfaufg}
\pfgruppe{}
\begin{pfaufg}{SH ’25}{31.}{}
Die Punkte B($-$0,5 | $-$1,05) und C(0,5 | 1,05) liegen auf einer Kurve. Bestimme die Gleichung der Geraden k durch B und C.
\rechenraster{2}
\fusshilfe{\mbox{31:~2,1; 2,1x}}
\end{pfaufg}
\begin{pfaufg}{}{32.}{}
Eine Taxifahrt von 4 km kostet 12{,}60 €. Eine Fahrt von 10 km kostet 25{,}20 €. Der Preis besteht aus einer Grundgebühr und einem festen Preis je km. Wie hoch ist die Grundgebühr?
\par\noindent \leerfeld[€] \par
\fusshilfe{\mbox{32:~$4{,}20$: Grundgebühr 4{,}20 €}}
\end{pfaufg}
\pfgruppe{}
\begin{pfaufg}{P10 ’25}{33.}{4 BE}
Der Graph einer linearen Funktion $f$ geht durch die Punkte $A(-2\,|\,6)$ und $B(3\,|\,-1{,}5)$. Zeichne den Graphen von $f$ in das Koordinatensystem. Entscheide, ob die Aussagen wahr oder falsch sind:
\par\smallskip\noindent \begin{tikzpicture}\begin{axis}[xmin=-4,xmax=4,ymin=-2,ymax=7,axis lines=middle,tick label style={font=\scriptsize},clip=true,/pgf/number format/use comma,axis line style={-{Stealth[length=2mm]}},every axis plot/.append style={line width=0.9pt},x=0.55cm,y=0.55cm,xtick distance=1,ytick distance=1,enlargelimits=false,grid=both,grid style={mbgitter,line width=0.3pt},major grid style={mbgitter},xlabel={$x$},ylabel={$y$},x label style={anchor=west},y label style={anchor=south}]\end{axis}\end{tikzpicture}
\par\smallskip\noindent \begingroup\renewcommand{\arraystretch}{1.5}\begin{tabular}{|p{8.0cm}|c|c|}\hline Aussage & wahr & falsch \\ \hline Der Graph von $f$ steigt monoton. & $\square$ & $\square$ \\ \hline Der Graph von $f$ schneidet die y-Achse in $(0\,|\,3)$ & $\square$ & $\square$ \\ \hline\end{tabular}\endgroup
\par\smallskip\noindent Gib eine Gleichung von $f$ an\par
\fusshilfe{\mbox{33:~f(x) = $-$1,5x + 3}}
\end{pfaufg}

\pfstufekopf[srechnerischangeradeu]{rechnerisch an Gerade und Parabel}{in 3 der letzten 5 Prüfungen}
\pfgruppe{}
\begin{pfaufg}{}{34.}{}
Setze $x = 4$ in den Term $3x - 2$ ein. Berechne den Wert.
\rechenraster{1}
\fusshilfe{\mbox{34:~$10$}}
\end{pfaufg}
\begin{pfaufg}{}{35.}{}
Kreuze an, ob der Graph von $f(x) = 4x - 1$ eine Gerade oder eine Parabel ist.
\pfkreuzzeile{Gerade}\pfkreuzzeile{Parabel}
\fusshilfe{\mbox{35:~Gerade}}
\end{pfaufg}
\begin{pfaufg}{}{36.}{}
Kreuze an, in welcher Form $f(x) = x^2 + 4x - 1$ geschrieben ist.
\pfkreuzzeile{Scheitelpunktform}\pfkreuzzeile{Normalform}
\fusshilfe{\mbox{36:~Normalform}}
\end{pfaufg}
\begin{pfaufg}{}{37.}{}
Du willst die Schnittpunkte der Parabel $f(x) = x^2 - 2x$ und der Geraden $g(x) = x + 4$ berechnen. Schreibe nur die Gleichung auf, mit der du anfängst.
\par\noindent \leerfeld[=]  \leerfeld \par
\fusshilfe{\mbox{37:~$x + 4$}}
\end{pfaufg}
\pfgruppe{}
\begin{pfaufg}{}{38.}{}
Setze $x = 4$ in $f(x) = 2x + 5$ ein. Berechne so den Funktionswert $f(4)$.
\par\noindent f(4) = \leerfeld \par
\fusshilfe{\mbox{38:~$13$}}
\end{pfaufg}
\begin{pfaufg}{}{39.}{}
Gib den Schnittpunkt der Parabel $f(x) = x^2 + 3x + 5$ mit der $y$-Achse an.
\par\noindent (\leerfeld[|\_\_)] \par
\fusshilfe{\mbox{39:~Schnittpunkt $(0|5)$}}
\end{pfaufg}
\begin{pfaufg}{}{40.}{}
Berechne die Nullstellen von $f(x) = (x - 1)^2 - 4$.
\par\noindent x\_1 = \leerfeld , x\_2 = \leerfeld \par
\fusshilfe{\mbox{40:~$x_1 = 3$, $x_2 = -1$}}
\end{pfaufg}
\begin{pfaufg}{NI ’26}{41.}{}
Prüfe mit einer Rechnung, ob P(10 | $-$1) auf dem Graphen von y$_2$ = $-$½x + 4 liegt.
\rechenraster{1}
\fusshilfe{\mbox{41:~$-$½ · 10 + 4 = $-$1 $\Rightarrow$ P liegt darauf}}
\end{pfaufg}
\begin{pfaufg}{NRW ’23}{42.}{}
Der Graph A hat die Gleichung y = $-$2x $-$ 3. Prüfe, ob P(4 | $-$11) auf ihm liegt.
\rechenraster{1}
\fusshilfe{\mbox{42:~$-$2 · 4 $-$ 3 = $-$11: P liegt auf A.}}
\end{pfaufg}
\begin{pfaufg}{VERA ’13}{43.}{}
Zeige: Der Punkt A(4 | 3) liegt auf der Geraden y = 1,5x $-$ 3.
\rechenraster{1}
\fusshilfe{\mbox{43:~1,5 · 4 $-$ 3 = 3 $\surd$}}
\end{pfaufg}
\begin{pfaufg}{}{44.}{}
Berechne $f(-6)$ für $f(x) = x^2$.
\par\noindent f(-6) = \leerfeld \par
\fusshilfe{\mbox{44:~$36$}}
\end{pfaufg}
\begin{pfaufg}{P10 ’17}{45.}{2 BE}
Die Gerade $g$ hat die Gleichung $y = \tfrac{1}{2}x + 1$. Auf ihr liegt der Punkt $A(-10\,|\,y)$. Berechne seine y-Koordinate.
\rechenraster{1}
\fusshilfe{\mbox{45:~y = $-$4}}
\end{pfaufg}
\pfgruppe{}
\begin{pfaufg}{P10 ’24}{46.}{2 BE}
Die Gerade $f$ ist der Graph von $f(x) = 4x + 1$. Die Parabel $p$ ist der Graph von $p(x) = x^2 - 4$. Prüfe mit einer Rechnung, ob der Punkt $T(-5\,|\,20)$ auf der Parabel $p$ liegt.
\rechenraster{1}
\fusshilfe{\mbox{46:~Nein, p($-$5) = 25 $-$ 4 = 21 $\neq$ 20}}
\end{pfaufg}
\pfsteckt{steckt auch in Nr. 72 (P10 ’23); Nr. 69 (P10 ’21)}
\pfstufekopf[sgleichungaufstellenu]{Gleichung aufstellen und rückwärts rechnen}{in 2 der letzten 5 Prüfungen}
\pfgruppe{}
\begin{pfaufg}{}{47.}{}
Ein Handwerker berechnet 45 € für die Anfahrt und 38 € je Arbeitsstunde. Unterstreiche den Anfangswert. Kreise die Änderung je Einheit ein.
\rechenraster{1}
\fusshilfe{\mbox{47:~Anfangswert 45 € , Änderung 38 € je Stunde}}
\end{pfaufg}
\begin{pfaufg}{}{48.}{}
Die Tabelle zeigt x-Werte und y-Werte. Jeder x-Wert wird mit demselben Faktor multipliziert. So entsteht der y-Wert. Wie groß ist der Faktor?
\par\smallskip\noindent \wertetabelle[3,6,9,12]{x}{y}{1,2,3,4}\par
\par\noindent Faktor: \leerfeld \par
\fusshilfe{\mbox{48:~3}}
\end{pfaufg}
\pfgruppe{}
\begin{pfaufg}{}{49.}{}
Ein Taxi kostet 4 € Grundgebühr und 2 € je km. Für x km zahlt man y Euro. Kreuze an, welche Gleichung passt. 
\pfkreuzzeile{$y = 4x + 2$}\pfkreuzzeile{$y = 6x$}\pfkreuzzeile{$y = 2x - 4$}\pfkreuzzeile{$y = 2x + 4$}
\fusshilfe{\mbox{49:~$2x + 4$}}
\end{pfaufg}
\pfgruppe{}
\begin{pfaufg}{}{50.}{}
Berechne die y-Werte zu $y = 3x$ für $x = 0, 1, 2$. Zeichne dann die Gerade durch diese Punkte in das Koordinatensystem.
\par\smallskip\noindent \begin{ksys}[xmin=-1,xmax=3,ymin=-1,ymax=7] \end{ksys}\par
\rechenraster{1}
\fusshilfe{\mbox{50:~Gerade durch $(0|0)$, $(1|3)$, $(2|6)$}}
\end{pfaufg}
\pfgruppe{}
\begin{pfaufg}{BY ’21}{51.}{}
Gib zur Wertetabelle einen passenden Term an. x: 0, 1, 2, 3, 4, 5; T(x): $-$7, $-$6, $-$5, $-$4, $-$3, $-$2.
\par\smallskip\noindent \sachtabelle{lrrrrrr}{x & 0 & 1 & 2 & 3 & 4 & 5}{T(x) & $-$7 & $-$6 & $-$5 & $-$4 & $-$3 & $-$2}\par
\fusshilfe{\mbox{51:~T(x) = x $-$ 7}}
\end{pfaufg}
\begin{pfaufg}{HH ’26}{52.}{}
Die Gerade g(x) = m · x + 3 geht durch P(2 | 7). Berechne die Steigung m.
\rechenraster{1}
\fusshilfe{\mbox{52:~7 = 2m + 3 $\Rightarrow$ m = 2}}
\end{pfaufg}
\begin{pfaufg}{NI ’24}{53.}{}
Ein Girokonto kostet im Tarif A 4,50 € im Monat und 0,15 € je Buchung. Stelle die Gleichung y = mx + b auf.
\par\smallskip\noindent \sachtabelle{lrr}{Tarif & Grundpreis & je Buchung}{}\par
\fusshilfe{\mbox{53:~y = 0,15x + 4,50}}
\end{pfaufg}
\begin{pfaufg}{NI ’22}{54.}{}
Ein E-Scooter kostet 1 € Freischaltung und 0,20 € je Minute. Stelle die Gleichung y = mx + b auf.
\rechenraster{1}
\fusshilfe{\mbox{54:~y = 0,20x + 1}}
\end{pfaufg}
\begin{pfaufg}{}{55.}{}
Denke dir eine Situation zur Gleichung $y = -15x + 300$ aus. Beschreibe sie und sage, was x und y bedeuten.
\rechenraster{1}
\fusshilfe{\mbox{55:~x Minuten, y Liter}}
\end{pfaufg}
\pfgruppe{}
\begin{pfaufg}{NI ’23}{56.}{}
Eine Ferienwohnung kostet 90 € je Nacht und 110 € Endreinigung. Stelle die Gleichung y = mx + b für den Gesamtpreis auf.
\rechenraster{1}
\fusshilfe{\mbox{56:~y = 90x + 110}}
\end{pfaufg}
\begin{pfaufg}{}{57.}{}
Ein Kino-Abo kostet einmalig 10 € und dazu 6 € je Film. Für x Filme zahlt man y Euro. Stelle eine Gleichung für y auf.
\par\noindent y = \leerfeld \par
\fusshilfe{\mbox{57:~$6x + 10$}}
\end{pfaufg}
\pfgruppe{}
\begin{pfaufg}{}{58.}{}
Ein Handwerker berechnet 35 € Anfahrt und 48 € je Stunde. Nele soll dazu eine Gleichung aufstellen. Sie schreibt $y = 35x + 48$. Finde den Fehler und gib die richtige Gleichung an.
\rechenraster{1}
\fusshilfe{\mbox{58:~$48x + 35$}}
\end{pfaufg}
\pfgruppe{}
\begin{pfaufg}{BY ’24}{59.}{}
Zu einem Term gibt es diese Wertetabelle: x: $-$1, 0, 1, 2; T(x): 2, 4, 6, 8. Kreuze den passenden Term an: 2x²; 6x; x · x + 4; x $-$ (x + 4); x + (x + 4).
\par\smallskip\noindent \sachtabelle{lrrrr}{x & $-$1 & 0 & 1 & 2}{T(x) & 2 & 4 & 6 & 8}\par
\fusshilfe{\mbox{59:~x + (x + 4) = 2x + 4}}
\end{pfaufg}
\pfgruppe{}
\begin{pfaufg}{P10 ’16}{\llap{\textasteriskcentered\,}60.}{2 BE}
Die 10b fährt mit dem Bus in eine Jugendherberge. Zwei Busfirmen machen ein Angebot: „Sonnenschein“ verlangt 200 € Grundpreis und 1,50 € je gefahrenem Kilometer, „Reiselust“ keinen Grundpreis und 2,00 € je Kilometer. Die Grafik zeigt für beide Angebote die Fahrtkosten in Abhängigkeit von der Strecke. Gib die Gleichung der Funktion an, deren Graph II ist (x: Strecke in km, y: Kosten in €).
\par\smallskip\noindent \sachtabelle{lrr}{Busfirma & Grundpreis & Preis je km}{Sonnenschein & 200 € & 1,50 €\\ Reiselust & – & 2,00 €}
\par\smallskip\noindent \begin{tikzpicture}\begin{axis}[xmin=0,xmax=450,ymin=0,ymax=800,tick label style={font=\scriptsize},clip=true,/pgf/number format/use comma,axis line style={-{Stealth[length=2mm]}},every axis plot/.append style={line width=0.9pt},width=11.5cm,height=7.5cm,grid=both,grid style={mbgitter,line width=0.3pt},minor tick num=1,axis lines=left,xlabel={Fahrtstrecke in km},ylabel={Fahrtkosten in €},label style={font=\small},scaled ticks=false,/pgf/number format/1000 sep={\,}]\addplot[black,domain=0:450,samples=120] {2*x};\node[font=\small,fill=white,inner sep=1pt,anchor=south west] at (axis cs:344.7,689.4) {$I$};\addplot[black,domain=0:450,samples=120] {1.5*x+200};\node[font=\small,fill=white,inner sep=1pt,anchor=south west] at (axis cs:285.3,627.95) {$II$};\end{axis}\end{tikzpicture}\par
\fusshilfe{\mbox{60:~y = 1,5x + 200}}
\end{pfaufg}
\begin{pfaufg}{P10 ’21}{61.}{2 BE}
Ordne jedem Sachverhalt die passende Gleichung zu. Sachverhalt 1: Eine Taxifahrt kostet 5 € Grundgebühr und 20 Cent je Kilometer. Sachverhalt 2: Ein Fitnessstudio verlangt im Monat 20 € Grundgebühr und 5 € für jeden gebuchten Kurs. Gleichungen: $y = 20x + 5$, $y = 5x + 20$, $y = 0{,}20x + 5$.
\par\smallskip\noindent ANKREUZ{|$y = 20x + 5$|$y = 5x + 20$|$y = 0{,}20x + 5$}{Sachverhalt 1|Sachverhalt 2}{fest}\par
\rechenraster{1}
\fusshilfe{\mbox{61:~0,20x + 5; 5x + 20}}
\end{pfaufg}
\begin{pfaufg}{P10 ’21}{\llap{\textasteriskcentered\,}62.}{3 BE}
Eine 40 cm hohe Kerze wird angezündet. Die Tabelle hält fest, wie hoch sie nach einigen Minuten noch ist. Das Abbrennen der Kerze beschreibt die Gleichung $y = -0{,}2x + 40$. Schreib zu $y$, $x$ und $40$, was sie hier bedeuten. Beispiel: $-0{,}2$ heißt, dass die Kerze in jeder Minute um 0,2 cm kürzer wird.
\par\smallskip\noindent \par\medskip\begingroup\centering\renewcommand{\arraystretch}{1.25}\begin{tabular}{|p{4.0cm}|p{8.5cm}|}\hline Teil der Gleichung & Bedeutung \\ \hline y & \leerzelle\\ \hline $-$0,2 & In jeder Minute wird die Kerze 0,2 cm kürzer.\\ \hline x & \leerzelle\\ \hline 40 & \leerzelle \\ \hline\end{tabular}\par\endgroup\medskip\par
\fusshilfe{\mbox{62:~y: Höhe in cm; x: Zeit in min; 40: Anfangshöhe in cm}}
\end{pfaufg}
\pfgruppe{}
\begin{pfaufg}{HH ’26}{63.}{}
Bestimme x durch Umformen: 4x $-$ 10 = $-$2x + 5.
\rechenraster{2}
\fusshilfe{\mbox{63:~6x = 15 $\Rightarrow$ x = 2,5}}
\end{pfaufg}
\pfgruppe{}
\begin{pfaufg}{P10 ’21}{64.}{2 BE}
Eine 40 cm hohe Kerze wird angezündet. Die Tabelle hält fest, wie hoch sie nach einigen Minuten noch ist. Die Gleichung $y = -0{,}2x + 40$ beschreibt das Abbrennen der Kerze ($x$: Zeit in Minuten, $y$: Höhe in cm). Bestimme, nach wie vielen Minuten die Kerze ganz heruntergebrannt ist.
\par\smallskip\noindent \sachtabelle{lccccccc}{Zeit in Minuten & 0 & 10 & 20 & 30 & 40 & … & 80}{Höhe der Kerze in cm & \leerzelle & 38 & 36 & 34 & 32 & … & \leerzelle}\par
\pfzwfrage{Gib an, welchen Wert $y$ hat, wenn die Kerze ganz abgebrannt ist.}
\fusshilfe{\mbox{64:~200 min (3 h 20 min)}}
\end{pfaufg}
\pfgruppe{}
\begin{pfaufg}{P10 ’22}{\llap{\textasteriskcentered\,}65.}{4 BE}
Maxi und Paula haben zum 16. Geburtstag Geld bekommen und legen jeden Monat etwas dazu, um den Führerschein zu bezahlen. Maxi beginnt mit 1 472 € und spart 22 € im Monat. Paula beginnt mit 990 € und spart 55 € im Monat. Stelle eine Gleichung für Paulas Guthaben auf ($y$: Guthaben in €, $x$: Zahl der Monate). Für Fahrschule und Prüfung braucht Paula 2 000 €. Bestimme, wie viele Monate sie mindestens sparen muss.
\par\smallskip\noindent \sachtabelle{lrr}{Name & Anfangsguthaben & monatlich}{Maxi & 1 472 € & 22 €\\ Paula & 990 € & 55 €}\par
\pfzwfrage{Setze $y = 2000$ in deine Gleichung ein und löse nach $x$ auf.}
\fusshilfe{\mbox{65:~y = 55x + 990; 19 Monate}}
\end{pfaufg}
\pfsteckt{steckt auch in Nr. 23 (P10 ’22); Nr. 91 (P10 ’23)}
\pfstufekopf[szeichnenundaussagenp]{zeichnen und Aussagen prüfen}{in 2 der letzten 5 Prüfungen}
\pfgruppe{}
\begin{pfaufg}{}{66.}{}
Eine Gerade hat die Gleichung $f(x) = 6$. Gib ihre Steigung m an und beschreibe, wie sie verläuft.
\par\noindent m = \leerfeld \par
\fusshilfe{\mbox{66:~$0$}}
\end{pfaufg}
\pfgruppe{}
\begin{pfaufg}{}{67.}{}
Welche Gerade ist steiler: $f(x) = 3x - 7$ oder $g(x) = 5x + 4$? Begründe ohne Rechnung.
\rechenraster{1}
\end{pfaufg}
\begin{pfaufg}{}{68.}{}
Schneiden sich die Geraden $f(x) = 3x - 7$ und $g(x) = 3x + 5$? Begründe ohne Rechnung.
\rechenraster{1}
\end{pfaufg}
\pfgruppe{}
\begin{pfaufg}{P10 ’21}{69.}{2 BE}
Die lineare Funktion $f$ hat die Gleichung $y = 3x + 1$. Kreuze die beiden Aussagen an, die für $f$ stimmen:
\pfkreuzzeile{Der Punkt $Q(5\,|\,0)$ liegt auf dem Graphen.}\pfkreuzzeile{Der Graph steigt monoton.}\pfkreuzzeile{Der Graph trifft die y-Achse in $P(0\,|\,1)$.}\pfkreuzzeile{Der Graph geht durch den Ursprung}
\pfzwfrage{Berechne $f(5)$.}
\fusshilfe{\mbox{69:~Aussage 2 und 3 richtig}}
\end{pfaufg}
\pfgruppe{}
\begin{pfaufg}{P10 ’23}{70.}{5 BE}
Im Koordinatensystem ist eine nach unten geöffnete Normalparabel $p$ gezeichnet. Die Gerade $f$ hat die Gleichung $f(x) = 4x - 2$. Zeichne $f$ in das Koordinatensystem. Gib die Koordinaten eines Punktes an, in dem sich die beiden Graphen schneiden. Entscheide, ob die Aussagen wahr oder falsch sind:
\par\smallskip\noindent \begin{tikzpicture}\begin{axis}[xmin=-5,xmax=5,ymin=-4,ymax=6,axis lines=middle,tick label style={font=\scriptsize},clip=true,/pgf/number format/use comma,axis line style={-{Stealth[length=2mm]}},every axis plot/.append style={line width=0.9pt},x=0.55cm,y=0.55cm,xtick distance=1,ytick distance=1,enlargelimits=false,grid=both,grid style={mbgitter,line width=0.3pt},major grid style={mbgitter},xlabel={$x$},ylabel={$y$},x label style={anchor=west},y label style={anchor=south}]\addplot[black,domain=-5:5,samples=120] {-((x+1)^2)+6};\node[font=\small,fill=white,inner sep=1pt,anchor=south west] at (axis cs:2,-3) {$p$};\end{axis}\end{tikzpicture}
\par\smallskip\noindent \begingroup\renewcommand{\arraystretch}{1.5}\begin{tabular}{|p{8.0cm}|c|c|}\hline Aussage & wahr & falsch \\ \hline Die Steigung von $f$ ist negativ. & $\square$ & $\square$ \\ \hline Gerade und Parabel haben zwei Schnittpunkte. & $\square$ & $\square$ \\ \hline $f$ schneidet die y-Achse in $(-2\,|\,0)$ & $\square$ & $\square$ \\ \hline\end{tabular}\endgroup\par
\pfzwfrage{Lies Steigung und y-Achsenabschnitt aus der Gleichung von $f$ ab.}
\fusshilfe{\mbox{70:~Gerade durch (0|$-$2) und (1|2); S(1|2)}}
\end{pfaufg}
\pfsteckt{steckt auch in Nr. 33 (P10 ’25)}
\pfstufekopf[sankreuzen]{ankreuzen}{in 1 der letzten 5 Prüfungen}
\pfgruppe{}
\begin{pfaufg}{}{71.}{}
Prüfe, ob der Punkt $P(3|8)$ auf der Geraden $f(x) = 3x - 1$ liegt.
\rechenraster{1}
\fusshilfe{\mbox{71:~Ja: $f(3) = 8$}}
\end{pfaufg}
\pfgruppe{}
\begin{pfaufg}{P10 ’23}{72.}{1 BE}
Die Gerade $f$ hat die Gleichung $f(x) = -7x + 3$. Einer dieser Punkte liegt nicht auf $f$. Kreuze ihn an:
\pfkreuzzeile{$(-1\,|\,10)$}\pfkreuzzeile{$(3\,|\,-24)$}\pfkreuzzeile{$(-4\,|\,31)$}
\fusshilfe{\mbox{72:~(3|$-$24), denn f(3) = $-$18 $\neq$ $-$24}}
\end{pfaufg}
\pfstufekopf[sendwertberechnen]{Endwert berechnen}{in 1 der letzten 5 Prüfungen}
\pfgruppe{}
\begin{pfaufg}{NI ’22}{73.}{}
Ein E-Scooter kostet 1 € fürs Freischalten und 0,20 € je Minute. Nina fährt 10 Minuten. Berechne die Kosten.
\rechenraster{1}
\fusshilfe{\mbox{73:~1 € + 10 · 0,20 € = 3 €}}
\end{pfaufg}
\begin{pfaufg}{}{74.}{}
Tom hat 85 € gespart. Jede Woche kommen 15 € dazu. Wie viel Geld hat Tom nach 12 Wochen?
\par\noindent \leerfeld[€] \par
\fusshilfe{\mbox{74:~$265$ €}}
\end{pfaufg}
\pfgruppe{}
\begin{pfaufg}{BY ’24}{75.}{}
Der Reaktionsweg eines Autos beträgt r(x) = 3 · x/10 Meter, wenn x die Geschwindigkeit in km/h ist. Wie lang ist der Reaktionsweg bei 130 km/h?
\fusshilfe{\mbox{75:~r(130) = 3 · 13 = 39 m}}
\end{pfaufg}
\begin{pfaufg}{NI ’23}{76.}{}
Ein Boiler erwärmt Wasser von 20 °C. Die Temperatur steigt pro Minute um 15 °C. Berechne die Temperatur nach 3 Minuten.
\rechenraster{1}
\fusshilfe{\mbox{76:~20 °C + 3 · 15 °C = 65 °C}}
\end{pfaufg}
\begin{pfaufg}{NI ’24}{77.}{}
Pauls Großeltern geben ihm sofort 100 € und dann jeden Monat 50 € dazu. Paul sagt: „Nach vier Monaten habe ich 300 €.“ Zeige, dass er recht hat.
\rechenraster{1}
\fusshilfe{\mbox{77:~100 € + 4 · 50 € = 300 €}}
\end{pfaufg}
\pfgruppe{}
\begin{pfaufg}{}{78.}{}
Ein Handwerker berechnet 45 € für die Anfahrt und 52 € je Stunde. Vorher hat er gesagt: Die Arbeit kostet höchstens 300 €. Reichen die 300 € für 5 Stunden Arbeit?
\rechenraster{1}
\fusshilfe{\mbox{78:~Nein: $45 + 5 \cdot 52 = 305$ €, das sind 5 € zu viel}}
\end{pfaufg}
\pfgruppe{}
\begin{pfaufg}{P10 ’21}{79.}{2 BE}
Eine 40 cm hohe Kerze wird angezündet. Die Tabelle hält fest, wie hoch sie nach einigen Minuten noch ist. Trag die zwei fehlenden Werte in die Tabelle ein.
\par\smallskip\noindent \sachtabelle{lccccccc}{Zeit in Minuten & 0 & 10 & 20 & 30 & 40 & … & 80}{Höhe der Kerze in cm & \leerzelle & 38 & 36 & 34 & 32 & … & \leerzelle}\par
\pfzwfrage{Gib an, um wie viele Zentimeter die Kerze in 10 Minuten kürzer wird.}
\fusshilfe{\mbox{79:~40 cm; 24 cm}}
\end{pfaufg}
\pfgruppe{}
\begin{pfaufg}{P10 ’22}{80.}{1 BE}
Maxi und Paula haben zum 16. Geburtstag Geld bekommen und legen jeden Monat etwas dazu, um den Führerschein zu bezahlen. Maxi beginnt mit 1 472 € und spart 22 € im Monat. Paula beginnt mit 990 € und spart 55 € im Monat. Nach einem Jahr hat Paula 1 650 € gespart. Berechne, wie viel Geld Maxi nach einem Jahr hat.
\par\smallskip\noindent \sachtabelle{lrr}{Name & Anfangsguthaben & monatlich}{Maxi & 1 472 € & 22 €\\ Paula & 990 € & 55 €}\par
\pfzwfrage{Berechne, wie viel Maxi in zwölf Monaten dazulegt.}
\fusshilfe{\mbox{80:~1736 €}}
\end{pfaufg}
\pfstufekopf[sgraphzumtarifzuordne]{Graph zum Tarif zuordnen}{in 1 der letzten 5 Prüfungen}
\pfgruppe{}
\begin{pfaufg}{}{81.}{}
Ein Tarif kostet 3 € je Kilogramm ohne Grundpreis. Das Koordinatensystem zeigt die Graphen A, B und C. Gib den Graphen an, der zum Tarif passt.
\par\smallskip\noindent \begin{ksys}[xmin=0,xmax=8,ymin=0,ymax=40,xstep=1,ystep=5,ablesen] \gerade{3}{1}{A} \gerade{3}{0}{B} \gerade{1}{3}{C} \end{ksys}\par
\par\noindent Graph \leerfeld \par
\fusshilfe{\mbox{81:~B}}
\end{pfaufg}
\pfgruppe{}
\begin{pfaufg}{P10 ’23}{82.}{2 BE}
Herr Mert braucht für seinen Umzug einen Transporter. Angebot 1: 27 € Miete je Tag, 350 km insgesamt frei, jeder weitere Kilometer kostet 0,36 €. Angebot 2: einmalig 120 € Miete für eine Woche, jeder Kilometer kostet 0,09 €. Die Kosten für einen Miettag sollen für beide Angebote in einem Diagramm dargestellt werden. Gib zu jedem Angebot das passende Diagramm an.
\par\smallskip\noindent \begin{tabular}[t]{@{}c@{}}{\small A}\enspace$\square$\\\begin{tikzpicture}\begin{axis}[width=3.9cm,height=3.4cm,scale only axis=false,axis lines=middle,xtick=\empty,ytick=\empty,xmin=0,xmax=1.08,ymin=0,ymax=1.08,clip=true,axis line style={-{Stealth[length=1.5mm]}},xlabel={\tiny Fahrstrecke in km},ylabel={\tiny Kosten in €},x label style={anchor=north east,at={(1,0)},yshift=-2pt},y label style={anchor=south west,at={(0,1)},xshift=1pt}]\addplot[black,line width=0.9pt,domain=0:1,samples=80] {0.85*x};\end{axis}\end{tikzpicture}\end{tabular}\hspace{0.4cm}\begin{tabular}[t]{@{}c@{}}{\small B}\enspace$\square$\\\begin{tikzpicture}\begin{axis}[width=3.9cm,height=3.4cm,scale only axis=false,axis lines=middle,xtick=\empty,ytick=\empty,xmin=0,xmax=1.08,ymin=0,ymax=1.08,clip=true,axis line style={-{Stealth[length=1.5mm]}},xlabel={\tiny Fahrstrecke in km},ylabel={\tiny Kosten in €},x label style={anchor=north east,at={(1,0)},yshift=-2pt},y label style={anchor=south west,at={(0,1)},xshift=1pt}]\addplot[black,line width=0.9pt,domain=0:1,samples=80] {0.3+0.35*x};\end{axis}\end{tikzpicture}\end{tabular}\hspace{0.4cm}\begin{tabular}[t]{@{}c@{}}{\small C}\enspace$\square$\\\begin{tikzpicture}\begin{axis}[width=3.9cm,height=3.4cm,scale only axis=false,axis lines=middle,xtick=\empty,ytick=\empty,xmin=0,xmax=1.08,ymin=0,ymax=1.08,clip=true,axis line style={-{Stealth[length=1.5mm]}},xlabel={\tiny Fahrstrecke in km},ylabel={\tiny Kosten in €},x label style={anchor=north east,at={(1,0)},yshift=-2pt},y label style={anchor=south west,at={(0,1)},xshift=1pt}]\addplot[black,line width=0.9pt] coordinates {(0,0.3) (0.4,0.3) (1,0.9)};\end{axis}\end{tikzpicture}\end{tabular}\hspace{0.4cm}\begin{tabular}[t]{@{}c@{}}{\small D}\enspace$\square$\\\begin{tikzpicture}\begin{axis}[width=3.9cm,height=3.4cm,scale only axis=false,axis lines=middle,xtick=\empty,ytick=\empty,xmin=0,xmax=1.08,ymin=0,ymax=1.08,clip=true,axis line style={-{Stealth[length=1.5mm]}},xlabel={\tiny Fahrstrecke in km},ylabel={\tiny Kosten in €},x label style={anchor=north east,at={(1,0)},yshift=-2pt},y label style={anchor=south west,at={(0,1)},xshift=1pt}]\addplot[black,line width=0.9pt] coordinates {(0,0.6) (0.4,0.6) (1,0.15)};\end{axis}\end{tikzpicture}\end{tabular}\par
\fusshilfe{\mbox{82:~Angebot 1: C; Angebot 2: B}}
\end{pfaufg}
\pfgruppe{}
\begin{pfaufg}{P10 ’16}{83.}{2 BE}
Die 10b fährt mit dem Bus in eine Jugendherberge. Zwei Busfirmen machen ein Angebot: „Sonnenschein“ verlangt 200 € Grundpreis und 1,50 € je gefahrenem Kilometer, „Reiselust“ keinen Grundpreis und 2,00 € je Kilometer. Die Grafik zeigt für beide Angebote die Fahrtkosten in Abhängigkeit von der Strecke. Gib zu jedem der beiden Graphen die passende Busfirma an. Begründe für „Sonnenschein“, warum dein Graph passt.
\par\smallskip\noindent \sachtabelle{lrr}{Busfirma & Grundpreis & Preis je km}{Sonnenschein & 200 € & 1,50 €\\ Reiselust & – & 2,00 €}
\par\smallskip\noindent \begin{tikzpicture}\begin{axis}[xmin=0,xmax=450,ymin=0,ymax=800,tick label style={font=\scriptsize},clip=true,/pgf/number format/use comma,axis line style={-{Stealth[length=2mm]}},every axis plot/.append style={line width=0.9pt},width=11.5cm,height=7.5cm,grid=both,grid style={mbgitter,line width=0.3pt},minor tick num=1,axis lines=left,xlabel={Fahrtstrecke in km},ylabel={Fahrtkosten in €},label style={font=\small},scaled ticks=false,/pgf/number format/1000 sep={\,}]\addplot[black,domain=0:450,samples=120] {2*x};\node[font=\small,fill=white,inner sep=1pt,anchor=south west] at (axis cs:344.7,689.4) {$I$};\addplot[black,domain=0:450,samples=120] {1.5*x+200};\node[font=\small,fill=white,inner sep=1pt,anchor=south west] at (axis cs:285.3,627.95) {$II$};\end{axis}\end{tikzpicture}\par
\fusshilfe{\mbox{83:~II beginnt bei 200 €}}
\end{pfaufg}
\pfstufekopf[starifevergleichen]{Tarife vergleichen}{in 1 der letzten 5 Prüfungen}
\pfgruppe{}
\begin{pfaufg}{}{84.}{}
Studio A kostet 50 € Aufnahmegebühr und 20 € je Monat. Studio B kostet 28 € je Monat ohne Aufnahmegebühr. Welches Studio ist für ein Jahr günstiger?
\rechenraster{1}
\fusshilfe{\mbox{84:~A ist günstiger}}
\end{pfaufg}
\pfgruppe{}
\begin{pfaufg}{}{85.}{}
Die Geraden von zwei Tarifen schneiden sich bei 200 km. Tim sagt: „Ab 200 km ist der Tarif mit Grundgebühr günstiger.“ Begründe, ob Tim recht hat.
\rechenraster{1}
\end{pfaufg}
\pfgruppe{}
\begin{pfaufg}{nach BW ’21}{86.}{}
Flyer drucken: Angebot 1 kostet 45 € fest und 0,39 € je Flyer. Angebot 2 kostet nur 0,45 € je Flyer. Welches Angebot ist bei 1 000 Flyern billiger?
\fusshilfe{\mbox{86:~Angebot 1: 435 € < 450 €}}
\end{pfaufg}
\begin{pfaufg}{HH ’26}{87.}{}
Tarif „Flat“: 25 € im Monat. Tarif „Easy“: 3 € Grundgebühr und 2 Cent pro Megabyte. Ermittle, ab wie viel MB „Flat“ günstiger ist.
\rechenraster{2}
\fusshilfe{\mbox{87:~x = 1 100 MB; ab mehr als 1 100 MB}}
\end{pfaufg}
\begin{pfaufg}{NI ’24}{88.}{}
Tarif A kostet 4,50 € im Monat und 0,15 € je Buchung. Tarif C kostet pauschal 7,90 €. Frau Weimann bucht im Monat etwa 25-mal. Welcher Tarif ist günstiger?
\rechenraster{2}
\fusshilfe{\mbox{88:~8,25 € > 7,90 € $\Rightarrow$ Tarif C}}
\end{pfaufg}
\begin{pfaufg}{VERA ’17}{89.}{}
Hin- und Rückfahrt Hamburg–Berlin kosten normal 140,00 €. Frau Schnell kauft für 230,00 € eine Bahncard 50. Damit zahlt sie jede Fahrkarte zum halben Preis. Wie oft muss sie hin- und zurückfahren, damit sich die Karte lohnt?
\rechenraster{3}
\fusshilfe{\mbox{89:~Mit Karte: 230 + 70n; mindestens 4-mal}}
\end{pfaufg}
\pfgruppe{}
\begin{pfaufg}{P10 ’16}{90.}{3 BE}
Die 10b fährt mit dem Bus in eine Jugendherberge. Zwei Busfirmen machen ein Angebot: „Sonnenschein“ verlangt 200 € Grundpreis und 1,50 € je gefahrenem Kilometer, „Reiselust“ keinen Grundpreis und 2,00 € je Kilometer. Die Grafik zeigt für beide Angebote die Fahrtkosten in Abhängigkeit von der Strecke. Die Jugendherberge liegt 175 km entfernt; der Bus fährt hin und zurück. Gib an, welche Busfirma dafür billiger ist. Am Ort kommt noch ein Tagesausflug von 100 km hinzu. Untersuche, ob dieselbe Firma dann noch die billigere ist, und begründe.
\par\smallskip\noindent \sachtabelle{lrr}{Busfirma & Grundpreis & Preis je km}{Sonnenschein & 200 € & 1,50 €\\ Reiselust & – & 2,00 €}
\par\smallskip\noindent \begin{tikzpicture}\begin{axis}[xmin=0,xmax=450,ymin=0,ymax=800,tick label style={font=\scriptsize},clip=true,/pgf/number format/use comma,axis line style={-{Stealth[length=2mm]}},every axis plot/.append style={line width=0.9pt},width=11.5cm,height=7.5cm,grid=both,grid style={mbgitter,line width=0.3pt},minor tick num=1,axis lines=left,xlabel={Fahrtstrecke in km},ylabel={Fahrtkosten in €},label style={font=\small},scaled ticks=false,/pgf/number format/1000 sep={\,}]\addplot[black,domain=0:450,samples=120] {2*x};\node[font=\small,fill=white,inner sep=1pt,anchor=south west] at (axis cs:344.7,689.4) {$I$};\addplot[black,domain=0:450,samples=120] {1.5*x+200};\node[font=\small,fill=white,inner sep=1pt,anchor=south west] at (axis cs:285.3,627.95) {$II$};\end{axis}\end{tikzpicture}\par
\pfzwfrage{Berechne die gesamte Strecke für Hin- und Rückfahrt.}
\pfzwfrage{Berechne die Kosten beider Firmen für diese Strecke.}
\pfzwfrage{Berechne die Kosten beider Firmen für die Strecke mit Tagesausflug.}
\fusshilfe{\mbox{90:~Reiselust (700 € < 725 €)}}
\end{pfaufg}
\pfgruppe{}
\begin{pfaufg}{P10 ’23}{\llap{\textasteriskcentered\,}91.}{5 BE}
Herr Mert braucht für seinen Umzug einen Transporter. Angebot 1: 27 € Miete je Tag, 350 km insgesamt frei, jeder weitere Kilometer kostet 0,36 €. Angebot 2: einmalig 120 € Miete für eine Woche, jeder Kilometer kostet 0,09 €. Herr Mert mietet den Transporter von Montag bis Freitag und fährt insgesamt 470 km. Berechne die Kosten für beide Angebote. Gib an, welches billiger ist. Stelle für Angebot 2 eine Gleichung auf, mit der man die Kosten einer Woche für jede Kilometerzahl berechnen kann.
\pfzwfrage{Berechne die Miete für fünf Tage bei Angebot 1.}
\rechenraster{4}
\fusshilfe{\mbox{91:~K(x) = 0,09x + 120}}
\end{pfaufg}

\pfpruefstein{P10 ’26}
\begin{pfaufg}{}{}{}
Die verschobene Normalparabel $p$ ist der Graph von $p(x) = (x-2)^2 - 2$ und ist im Koordinatensystem gezeichnet.
\par\smallskip\noindent \begin{tikzpicture}\begin{axis}[xmin=-2,xmax=5,ymin=-2,ymax=5,axis lines=middle,tick label style={font=\scriptsize},clip=true,/pgf/number format/use comma,axis line style={-{Stealth[length=2mm]}},every axis plot/.append style={line width=0.9pt},x=0.55cm,y=0.55cm,xtick distance=1,ytick distance=1,enlargelimits=false,grid=both,grid style={mbgitter,line width=0.3pt},major grid style={mbgitter},xlabel={$x$},ylabel={$y$},x label style={anchor=west},y label style={anchor=south}]\addplot[black,domain=-2:5,samples=120] {((x-2)^2)-2};\node[font=\small,fill=white,inner sep=1pt,anchor=south west] at (axis cs:4.44,3.9536) {$p$};\end{axis}\end{tikzpicture}\par
\end{pfaufg}
\begin{pfaufg}{}{*a)}{2 BE}
Die lineare Funktion $f$ hat die Gleichung $f(x) = -2x + 2$. Zeichne ihren Graphen in das Koordinatensystem.
\rechenplatz{3}
\fusshilfe{\mbox{Prüfstein a):~Gerade durch (0|2) und (1|0), fallend}}
\end{pfaufg}
\begin{pfaufg}{}{b)}{2 BE}
Prüfe mit einer Rechnung, ob der Punkt $P(-4\,|\,10)$ auf dem Graphen von $f(x) = -2x + 2$ liegt.
\rechenplatz{3}
\fusshilfe{\mbox{Prüfstein b):~Ja, f($-$4) = $-$2 · ($-$4) + 2 = 10}}
\end{pfaufg}
\begin{pfaufg}{}{c)}{1 BE}
Gib den Scheitelpunkt von $p$ an: $S(\ \ |\ \ )$.
\rechenplatz{3}
\fusshilfe{\mbox{Prüfstein c):~S(2|$-$2)}}
\end{pfaufg}
\begin{pfaufg}{}{*d)}{3 BE}
Eine zweite Parabel $q$ hat die Gleichung $q(x) = -(x-2)^2 - 2$. Gib an, wie viele Punkte $p$ und $q$ gemeinsam haben. Begründe ohne zu rechnen.
\rechenplatz{3}
\fusshilfe{\mbox{Prüfstein d):~genau ein gemeinsamer Punkt S(2|$-$2)}}
\end{pfaufg}
\end{document}